\documentclass[11pt]{article}
\usepackage[margin=1in]{geometry}
\usepackage{amsmath,amssymb,amsthm,hyperref}
\newtheorem{theorem}{Theorem}
\title{Characters of the free Lie algebra with compatible brackets}
\author{Wentao Li}
\date{November 20, 2026}
\begin{document}\maketitle
\begin{abstract}
We give a finite multinomial sum for the value of the multigraded character
on every conjugacy class of the symmetric group. It applies to any number of
compatible Lie brackets and yields the multiplicities of all Specht modules.
On a single cycle the sum reduces to ordered factorizations of the cycle
length. On a prime cycle the trace is $-1$ in every bracket multidegree.
\end{abstract}
Work in characteristic zero. For $\mu=(\mu_1,\ldots,\mu_k)$ of sum $n-1$,
$\operatorname{Lie}(\mu)$ is the component multilinear in $n$ distinct
generators in which bracket $j$ is used $\mu_j$ times. The brackets are
compatible, meaning every linear combination is a Lie bracket. Write
$\chi_\lambda(x)=\sum_\mu x^\mu
\operatorname{tr}(\sigma_\lambda\mid\operatorname{Lie}(\mu))$ for the
character at cycle type $\lambda=(1^{m_1}2^{m_2}\cdots)\vdash n$.
For any partition $\eta$ let $m_j(\eta)$ denote its multiplicity of $j$,
and $z_\eta=\prod_j j^{m_j(\eta)}m_j(\eta)!$.
Use $h_d$ for the complete homogeneous symmetric polynomial and
$x^i=(x_1^i,\ldots,x_k^i)$.

\begin{theorem}
For each pair $(i,\eta)$ with $i\ge1$, $|\eta|\ge2$ and $i|\eta|\le n$ put
\[
 w_{i,\eta}=\frac{(-1)^{\ell(\eta)}}{z_\eta}h_{|\eta|-1}(x^i).
\]
Sum over arrays of nonnegative integers $a_{i,\eta}$ satisfying
$\sum_{i,\eta}i(|\eta|-1)a_{i,\eta}=n-1$. Define
\[
 N_i=m_i+\sum_\eta a_{i,\eta},\qquad
 r_i=\mathbf1_{i=1}
       +\sum_{\substack{h\mid i\\h<i}}\sum_\eta a_{h,\eta}m_{i/h}(\eta).
\]
Keep only arrays for which $N_i=r_i+\sum_\eta a_{i,\eta}m_1(\eta)$ for
$1\le i\le n$, and $r_i>0$ whenever $N_i>0$. Then
\[
 \boxed{\displaystyle
 \chi_\lambda(x)=z_\lambda\sum_a
 \prod_{i:N_i>0}\left[
 \frac{r_i}{N_i}\binom{N_i}{m_i,(a_{i,\eta})_\eta}
             \prod_{\eta}w_{i,\eta}^{\,a_{i,\eta}}\right].}
                                                                    \tag{1}
\]
In particular, the sum is explicit in the cycle counts and the colors.
\end{theorem}

\begin{proof}
Use the Frobenius power sums $p_j$ for the symmetric-group variables.
By Theorem 1.8 of Gonz\'alez D'Le\'on \cite{g}, the multigraded Frobenius
series $T$ is the negative of the plethystic inverse of
$-\sum_{d\ge1}h_{d-1}(x)h_d$. Put $T_i=T(x^i;p_i,p_{2i},\ldots)$.
Using $h_d=\sum_{\eta\vdash d}p_\eta/z_\eta$ in this inverse identity gives
\[
 T_i=p_i+\sum_{|\eta|\ge2} w_{i,\eta}
                          \prod_{j\ge1}T_{ij}^{\,m_j(\eta)}.       \tag{2}
\]
In particular, $p_j[-T]=-T_j$, contributing one minus sign for each
part of $\eta$ prior to isolating the linear term.

Expand $T_1$ by rooted ordered trees. A leaf of type $i$ supplies $p_i$;
an internal node of type $i$ and signature $\eta$ has $m_j(\eta)$ child
slots of type $ij$ and supplies $w_{i,\eta}$. Fix the order of those slots
by increasing $j$, with distinct positions for repeats. This gives (2) by
root decomposition, with ordinary powers counting ordered choices. The total
type of the children exceeds the parent's type by $i(|\eta|-1)$.
Telescoping proves the array bound in the theorem and $i|\eta|\le n$ in
every contributing tree. The expansion through leaf mass $n$ is finite.

To count the trees with given leaf multiplicities and internal counts,
remove the edges on which type increases. For each $i$ what remains of
type $i$ is an ordered forest with $r_i$ roots, in the incoming slots from
smaller types (or the initial root for $i=1$). To see that the choices
multiply, construct in increasing order of $i$, ordering the available
incoming slots by their positions in the partial plane tree. For each
same-type forest, the prescribed signatures uniquely reinsert all its
slots to larger types. This construction reverses cutting.

For fixed type $i$, let $b$ mark a leaf and $c_\eta$ an internal signature.
Its plane trees are generated by $Y=t(b+\sum_\eta c_\eta
Y^{m_1(\eta)})$. By ordinary Lagrange inversion,
\[
 [t^{N_i}]Y^{r_i}=\frac{r_i}{N_i}[s^{N_i-r_i}]
               (b+\sum_\eta c_\eta s^{m_1(\eta)})^{N_i}.
\]
The coefficient of $b^{m_i}\prod c_\eta^{a_{i,\eta}}$ is exactly the
multinomial factor in (1), with the stated balance condition. Empty
types contribute one. Multiplying the weights and using the Frobenius
identity $\chi_\lambda(x)=z_\lambda[p_\lambda]T_1$ proves (1).
\end{proof}

\paragraph{Irreducible multiplicities.}
For $\rho\vdash n$, take $l\ge\ell(\rho)$, pad $\rho$ with zeroes, and put
$\delta=(l-1,\ldots,0)$, $\Delta(y)=\prod_{a<b}(y_a-y_b)$.
The multiplicity of $S^\rho$ in $\operatorname{Lie}(\mu)$ is
\[
 [x^\mu y^{\rho+\delta}]\Delta(y)
   \sum_{\lambda\vdash n}\frac{\chi_\lambda(x)}{z_\lambda}
     \prod_{j=1}^n (y_1^j+\cdots+y_l^j)^{m_j}.
\]
Indeed the sum is the Frobenius characteristic in ordinary variables.
Multiplication by the Vandermonde sends each Schur polynomial to its
alternant. Its strictly decreasing exponent vector detects precisely the
requested Schur coefficient. Formula (1) makes this expression a finite
sum involving only multinomials and symmetric polynomials of specified degree.

\paragraph{One cycle.}
For $\lambda=(n)$ the tree has exactly one leaf, so it is a chain. Internal
signatures are single parts $d\ge2$. Their product over the chain is $n$,
and the denominator $\prod d$ cancels $z_{(n)}=n$. Thus
\[
 \chi_{(n)}(x)=
 \sum_{\substack{d_1\cdots d_s=n\\d_j\ge2}}(-1)^s
     \prod_{j=1}^s h_{d_j-1}(x^{d_1\cdots d_{j-1}}).
\]
Include the empty factorization for $n=1$. For prime $n$ this is
$-h_{n-1}(x)$, giving trace $-1$ for each weak composition $\mu$ of $n-1$.
For just one bracket it gives the arithmetic M\"obius function at $n$,
as in the classical Lie character.

\paragraph{Status and credit.} This answers the character-value and
multiplicity requests of Question 7.13 in \cite{g} with a signed finite
sum. It gives no positive tableau rule and does not settle the separate
Schur-positivity Conjecture 7.14. The plethystic Frobenius formula is
Gonz\'alez D'Le\'on's. Lagrange inversion for cycle indices is classical
\cite{gl}, and Dotsenko and Khoroshkin gave the two-bracket character earlier
\cite{dk}. Research draft prepared with assistance from OpenAI Codex.
Independent proof and priority review is pending.
\begin{thebibliography}{9}
\bibitem{g} R. S. Gonz\'alez D'Le\'on, \emph{On the free Lie algebra with
multiple brackets}, Adv. Appl. Math. 79 (2016), 37--97.
\url{https://arxiv.org/abs/1408.5415}
\bibitem{gl} I. M. Gessel and G. Labelle, \emph{Lagrange inversion for
species}, J. Combin. Theory Ser. A 72 (1995), 95--117. \url{https://doi.org/10.1016/0097-3165(95)90030-6}
\bibitem{dk} V. Dotsenko and A. Khoroshkin, \emph{Character formulas for
the operad of two compatible brackets and for the bi-Hamiltonian operad},
Functional Analysis and Its Applications 41 (2007), 1--17. \url{https://doi.org/10.1007/s10688-007-0001-3}
\end{thebibliography}
\end{document}
